The \wikilink{a2s17--b010}{direct-sum and quotient construction of limits} makes the behavior of torsion explicit. In this subsection $R$ is a commutative integral domain with $1\ne0$. An element is torsion if it is killed by a nonzero scalar, as in \wikilink{a2s08--b001}{the torsion definition}.

\begin{lemma}\label{lem-torsion-sums-quotients}
A direct sum of torsion $R$-modules is torsion, and every quotient of a torsion module is torsion.
\end{lemma}
\begin{proof}\label{proof-torsion-sums-quotients}
An element $x=(x_i)$ of a direct sum has finite support. Choose for each nonzero coordinate a nonzero $a_i\in R$ with $a_ix_i=0$. Their finite product is nonzero because $R$ is a domain, and kills every coordinate because $R$ is commutative. If the support is empty, $1$ kills $x=0$. Hence the direct sum is torsion. An element of a quotient is the image of an element of the original module, so any nonzero scalar annihilating that lift also annihilates the image.
\end{proof}

\begin{proposition}\label{prop-torsion-colimits}
The colimit of any small diagram of torsion $R$-modules, computed in the category of $R$-modules, is torsion. In particular a directed direct limit of torsion modules is torsion.
\end{proposition}
\begin{proof}\label{proof-torsion-colimits}
For a diagram $(M_i)$, form $D=\bigoplus_iM_i$ with inclusions $\iota_i$. Let $N$ be the submodule generated by
\[
\iota_j(f_\alpha(m))-\iota_i(m),
\qquad \alpha:i\to j,\quad m\in M_i,
\]
where $f_\alpha$ is the map assigned to an arrow of the diagram. Compatible maps $M_i\to Y$ extend uniquely to a map $D\to Y$ and kill the indicated compatibility relations. Consequently $D/N$ has the universal property of the colimit. The lemma makes $D$ torsion and then makes its quotient $D/N$ torsion.

For a directed system, one can also represent an element at a single stage: a finite sum of stage representatives moves to a common later stage. Any nonzero scalar killing that representative kills its image in the direct limit.
\end{proof}

\begin{example}\label{ex-inverse-limit-torsion-failure}
Inverse limits need not be torsion even when every stage is torsion. Take $R=\mathbb Z$ and the system $\mathbb Z/p^n\mathbb Z$ with reduction maps. Its inverse limit is the \wikilink{s21--b009}{ring of p-adic integers}. The compatible tuple
\[
u=(1\bmod p^n)_{n\ge1}
\]
has infinite additive order. If a nonzero integer $a$ killed $u$, then $p^n$ would divide $a$ for every $n$, impossible once $p^n>|a|$. Thus the inverse limit is not a torsion $\mathbb Z$-module, although every stage is.
\end{example}

\begin{remark}\label{remark-uniform-annihilator}
If a single nonzero scalar kills every module in an inverse system, it also kills every compatible tuple, so that particular inverse limit is torsion. The example above has no such uniform nonzero annihilator. The failure of torsion preservation should not be confused with the separate question of exactness of inverse limits.
\end{remark}
